\chapter{Cycles and min-max families}\label{sec:minmax-cycles}
Pitts' monograph \cite{Pitts81} is the foundational source for discrete homotopies, almost minimizing varifolds, and regularity. We introduce the continuous notation used by Marques--Neves \cite{MNIndex}, Sections 1--2, and use Li's July 7, 2019 preliminary notes \cite{LiAP} as a guide to the discrete construction. A homotopy-class width and a volume-spectrum width are different definitions.

\begin{definition}[Mod-two boundary cycles and topologies]\label{def:minmax-cycles}
Let $(M^{n+1},g)$ be smooth, closed, and connected, with $n\geq1$. Write $\mathcal Z_n^0(M;\mathbb Z_2)$ for the flat component of zero in the space of integral mod-two $n$-cycles: these are boundaries $T=\partial U$ of mod-two $(n+1)$-chains. We use $\mathcal F$ for the flat metric, $\mathbf M$ for mass, and $\mathbf F_V$ for a metric inducing weak varifold convergence on bounded-mass sets. Set
\[
 \mathbf F(T,S)=\mathcal F(T-S)+\mathbf F_V(|T|,|S|).
\]
Here $|T|$ is the associated unoriented varifold. We take the usual bounded-Lipschitz varifold metric on the compact Grassmann bundle: its test functions have absolute value and Lipschitz constant at most one. In particular the constant test function one gives $|\|V\|(M)-\|W\|(M)|\leq\mathbf F_V(V,W)$. The flat topology can lose area by cancellation, while the varifold term records area and tangent planes. Source: \cite{MNIndex}, Section 2. The integer currents of the foundational chapter and these mod-two boundary cycles have different coefficient conventions.
\end{definition}

\begin{definition}[A fixed homotopy-class width]\label{def:homotopy-width}
Let $X$ be a finite simplicial complex and $\Phi:X\to\mathcal Z_n^0(M;\mathbb Z_2)$ be $\mathbf F$-continuous. Let $\Pi$ contain all $\mathbf F$-continuous maps from $X$ that are homotopic to $\Phi$ in the flat topology. Its width is
\[
 \mathbf L(\Pi)=\inf_{\Psi\in\Pi}\sup_{x\in X}\mathbf M(\Psi(x)).
\]
A minimizing sequence $\Phi_i\in\Pi$ has maximal masses converging to $\mathbf L(\Pi)$. Its critical set consists of varifold limits of slices $|\Phi_{i_j}(x_j)|$ whose masses converge to $\mathbf L(\Pi)$. This is the convention of \cite{MNIndex}, Section 1. The discrete homotopy formulation in \cite{Pitts81}, Section 4.1, requires a separate interpolation argument to compare the constructions.
\end{definition}

\begin{definition}[Mass fineness of a discrete family]\label{def:mass-fineness}
On the vertices of a finite cubical subdivision, a map $\phi$ to cycles has mass fineness
\[
 \mathbf f(\phi)=\max\{\mathbf M(\phi(x)-\phi(y)):x,y\text{ are adjacent vertices}\}.
\]
Almgren--Pitts sequences use successively finer subdivisions, uniformly bounded masses, and fineness tending to zero; their homotopies must satisfy corresponding fineness conditions. Small flat distances alone do not supply these conditions. Sources: \cite{Pitts81}, Section 4.1; \cite{LiAP}, Section 4.1; \cite{LMN}, Section 2.9.
\end{definition}

\begin{remark}[Smooth sweepouts in dimension three]\label{rem:smooth-sweepouts}
Colding--De Lellis \cite{CDL}, Section 1.1, formulate one-parameter families of surfaces on a closed three-manifold and allow specified degenerations at finitely many points. Their saturated classes are closed under parameter-dependent ambient isotopies. This provides a geometric route to the pull-tight and replacement arguments, alongside Pitts' cycle-space formulation. The parameter space, allowed degenerations, and saturation condition are part of that construction; they must be checked before transferring a width or a genus conclusion to another setting. The author's publication page also lists an erratum, which should be consulted when importing detailed proofs from this survey.
\end{remark}

\chapter{Almost minimizing varifolds and existence}\label{sec:almost-minimizing}
A critical limit needs more than stationarity. Pitts arranges that it resists small local mass-decreasing deformations in annuli; replacement arguments then connect this variational condition to regularity.

\begin{definition}[Almost minimizing in an open set]\label{def:almost-minimizing}
Fix an open $U\subset M$ and $\varepsilon,\delta>0$. A mod-two cycle $T$ is $(\varepsilon,\delta)$-almost minimizing in $U$ in the mass formulation if every finite sequence $T=T_0,\ldots,T_q$ of cycles with
\[
 \operatorname{spt}(T_i-T)\subset U,\quad
 \mathbf M(T_i-T_{i-1})\leq\delta,\quad
 \mathbf M(T_i)\leq\mathbf M(T)+\delta
\]
satisfies $\mathbf M(T_q)\geq\mathbf M(T)-\varepsilon$.
A varifold $V$ is almost minimizing in $U$ in this formulation if for every $\varepsilon>0$ some $\delta>0$ and such a $T$ satisfy $\mathbf F_V(V,|T|)<\varepsilon$.
The annular condition requires, for every $p\in M$, a radius $r_p>0$ such that this property holds in each $B_t(p)\setminus\overline B_s(p)$ with $0<s<t<r_p$. Compare \cite{LiAP}, Definition 2.1 and Theorem 2.2, and \cite{Pitts81}, Sections 3.1--3.9. The cited sources compare flat and mass formulations after restricting to smaller open sets; those equivalence proofs are not included here.
\end{definition}

\begin{theorem}[Pitts' original smooth existence range]\label{thm:pitts-existence}
Every smooth closed Riemannian manifold $M^{n+1}$ with $2\leq n\leq5$ contains a nonempty smooth closed embedded minimal hypersurface. This is the smooth-metric specialization of \cite{Pitts81}, Introduction, Theorem A, and Theorem 7.13. Pitts' proof is an external dependency. This statement records the dimension range in the supplied book; later regularity extensions are not being attributed to that original statement.
\end{theorem}

\begin{theorem}[Realization of a positive homotopy width]\label{thm:ap-realization}
Let $M^{n+1}$ be smooth, closed, and connected, with $2\leq n\leq6$. Let $\Pi$ and a minimizing sequence be as in Definition \ref{def:homotopy-width}, and suppose $\mathbf L(\Pi)>0$. The critical set contains an integral varifold
\[
 V=\sum_{j=1}^N m_j|\Sigma_j|,\qquad
 \|V\|(M)=\mathbf L(\Pi),
\]
where the $m_j$ are positive integers and the $\Sigma_j$ are pairwise disjoint smooth closed embedded minimal hypersurfaces. This modern formulation, including the Schoen--Simon regularity input, is stated in \cite{MNIndex}, Section 1, p. 4; see also \cite{LiAP}, Section 5. Its proof remains external. Neither connectedness nor multiplicity one is part of this conclusion.
\end{theorem}

\begin{remark}[The proof dependencies in Pitts' construction]\label{rem:pitts-proof-route}
The source route is: discrete homotopies and a nontrivial class (\cite{Pitts81}, Sections 4.1--4.7), pull-tight to obtain stationarity (Section 4.3), the combinatorial argument producing almost minimizing critical limits (Sections 4.8--4.13), and replacement and regularity arguments (Chapters 5--7). Li \cite{LiAP}, Sections 3--5, separates the same topological, variational, and analytic steps. Colding--De Lellis \cite{CDL} supplies the three-dimensional smooth-family perspective. A deformation that establishes stationarity alone does not complete the regularity step. The continuous pull-tight bridge is now proved in Theorem \ref{thm:pulled-tight-minimizing-sequence}, using the explicit deformation of Theorem \ref{thm:continuous-pulltight-homotopy}. Discrete interpolation, the combinatorial almost-minimizing selection, and replacement regularity remain external dependencies.
\end{remark}
